\chapter{RESULTADOS ESPERADOS}
\label{cap:resultados}

A partir dos objetivos estabelecidos, espera-se como resultado geral o desenvolvimento de um sistema completo para monitoramento e correção da execução de exercícios de força com halteres instrumentados com sensores inerciais. Esse sistema deverá ser capaz de fornecer feedback automático em tempo real ao usuário, além de incorporar informações provenientes de profissionais por meio do modo treinador, criando uma solução prática, acessível e cientificamente fundamentada. Assim, o trabalho busca avançar em relação às propostas existentes ao combinar portabilidade, baixo custo e personalização do aprendizado de máquina, com potencial de aplicação tanto em ambientes domésticos quanto em academias.  

Do ponto de vista técnico, espera-se a construção de um protótipo funcional de halteres instrumentados com IMUs, que permita a coleta precisa de dados de movimento em diferentes condições de execução. Esse protótipo será validado inicialmente em um exercício de força selecionado, a fim de controlar variáveis e estruturar o protocolo de coleta. Os dados obtidos deverão possibilitar a extração de características relevantes para a análise, bem como a implementação e avaliação de algoritmos de aprendizado de máquina voltados para a identificação da qualidade da execução. A integração desses algoritmos ao aplicativo móvel deverá permitir não apenas a retroalimentação em tempo real no modo usuário, mas também a incorporação das avaliações feitas no modo treinador, ampliando a base de dados e fortalecendo a capacidade de personalização do sistema.  

No âmbito acadêmico, espera-se que o trabalho contribua para a literatura ao propor uma abordagem de instrumentação de equipamentos de treino em substituição a soluções vestíveis ou de visão computacional. Pretende-se demonstrar a viabilidade do uso de sensores acoplados diretamente aos halteres como alternativa menos invasiva e mais prática, gerando evidências sobre sua acurácia e aplicabilidade. Além disso, busca-se oferecer uma discussão crítica sobre o uso de aprendizado de máquina em contextos de alta variabilidade individual, apontando a relevância de estratégias que incorporem feedback humano no treinamento dos modelos. Como resultado, espera-se a publicação de artigos em periódicos e congressos da área, contribuindo tanto para a comunidade científica quanto para aplicações práticas em saúde, esporte e tecnologia assistiva.  

Embora se preveja resultados positivos, algumas limitações são esperadas. A utilização inicial de apenas um exercício restringe a abrangência da validação e poderá limitar a generalização imediata do sistema. Diferenças individuais entre participantes, como experiência prévia, força e antropometria, também poderão afetar a consistência dos dados e a acurácia dos algoritmos. Além disso, o tempo e a complexidade da coleta de dados podem limitar o tamanho da amostra disponível para treinamento e validação, o que representa um desafio para métodos mais complexos de aprendizado profundo. Questões de usabilidade, como o peso adicional do sensor e a aceitação do protótipo pelos usuários, também podem influenciar os resultados. Ainda assim, tais limitações serão reconhecidas e discutidas, de forma a orientar melhorias futuras e indicar caminhos para a expansão da pesquisa.  


%\section{Descrição Geral}
%\section{Resultados Técnicos Esperados}
%\section{Resultados Acadêmicos Esperados}
%\section{Limitações Esperadas}